\subsection*{Shared external-input catalogue}
Eight explicit external predicates are implemented and uninstantiated: SS, TV, PNT, MW, MW-V, the standalone KM-U interface, KM-R and KM-Compatible. OR is a derived specification proved from TV, not an independent assumption about the author's other paper. Their actual names and types are in \texttt{ExternalInputs.lean}, \texttt{OrbitInputs.lean}, \texttt{KMTTInterfaces.lean} , \texttt{PrimeMassInputs.lean} and \texttt{LongVariation.lean}. The compatible predicate includes the same-parameter comparison and mass witnesses; the standalone KM-U predicate is not an extra premise when that compatible input is used. A catalogue entry is not a Lean axiom.

\ifDetailed
\textbf{SS: quoted polynomial-phase bound.} Shao--Sun, \emph{Linear equations in Piatetski--Shapiro primes}, arXiv:2605.18676v1, Theorem 2.1 and Definition 4.5. The source states the coefficient-uniform polynomial-phase consequence directly. \texttt{ShaoSunPhaseInput} quotes that consequence with a saving uniform throughout the degree-dependent interval of $\gamma$. Restricting real cutoffs to integers and extending $X\geq3$ to the two smaller positive cutoffs by enlarging a $\gamma$-dependent constant are included in the external translation. The general nilsequence/circle construction is optional unformalized exposition, not an additional required bridge for this quoted consequence. Both main branches depend on SS.

\textbf{TV: finite polynomial estimate.} Ter\"av\"ainen, \emph{Pointwise convergence of ergodic averages with M\"obius weight}, arXiv:2401.03174v2, Theorem 4.1. Its normalization is $M^{-d-1}$, its support scale is $M^d$, and its control is $u^{d+1}$. \texttt{FiniteVNInput} uses integer cutoffs, positive degrees and a universal bound on every phase correlation. The supremum/sign/monotonicity equivalence with the published norm is the unformalized formulation bridge TV-B. The constant depends on the fixed polynomial family. Both branches now depend on TV: it gives Proposition 2.3 and the proved orbit transference lemma. The same degree-only exponent is retained.

\textbf{KM-U and KM-R: the two KMTT finite theorems.} Krause--Mousavi--Tao--Ter\"av\"ainen, arXiv:2608.19525v1, Proposition 5.2, (5.9)--(5.12), gives weighted-to-unweighted discrepancy. \texttt{KMTTUnweightedInput} records its scale and density restrictions, positive discrepancy constant and retained parameter range. Lemma 5.3, (5.25)--(5.28), gives weak regularity, including its small-scale alternative. \texttt{KMTTRegularityInput} records coordinate resolutions $Q^{-O(1)}\varepsilon^{O(1)}M_0^{\nu_i}$, where $\nu_0=\deg P_1$ and $\nu_i=\deg P_i$ for $i\geq1$, and the common modulus $Q^b$ with $b=O_d(1)$. For fixed $C,d,\ell$, the regularity divisibility condition is $\mathcal L(\lceil\varepsilon^{-A'}\rceil)\mid Q$, with $0<\varepsilon<1/2$ and $1\leq M_0\leq(N/Q^{d-1})^{1/d}$. It applies to every tuple of $1$-bounded functions supported on $[N]$. Its alternatives are $M_0\leq C_{\rm sm}(Q/\varepsilon)^T$, or approximation of the actual unweighted counting form to error at most $\varepsilon$ by the respective atomwise conditional averages. Each atom is a progression whose cardinality lies between its resolution and twice that resolution. The positive implicit prefactors are retained explicitly. Compatibility uses (5.35)--(5.36). Integer exponent, real-observable, support and finite-form translations remain explicitly external as KM-U-B and KM-R-B; the local-factor density deduction and all scale bookkeeping are proved internally. Those original arguments are not included in an external input.

\textbf{KM-M and KM-J: compatible prime mass.} Source (5.2) gives mass in one fixed initial fraction, and (5.15) gives the full-mass lower bound, retaining possible exceptional-zero bias. \texttt{PrimeProgressionMassWitness} requires one fixed cutoff $0<\tau\leq\tau_{\max}$ and truncated mass at least $4c_1$ times the full mass. \texttt{KMTTCompatibleInput} jointly selects the comparison record and this witness, after a positive fixed $\tau_{\max}$ has been chosen. KM-J is the unformalized source-to-interface constant-selection bridge: choose the source's sufficiently small cutoff below $\tau_{\max}$, choose a smaller positive Lemma 5.1 constant $c_1$, and rerun Proposition 5.2's proof with enlarged fixed constants and thresholds. This is supported by its proof, not by monotonicity of the already fixed theorem statement. Separate existential comparison and mass inputs do not imply a compatible pair, and no mass witness is claimed for every arbitrary record. In the application $\tau_{\max}=1/(10C(d+1))$, preserving the fixed-parameter dependence. Only recurrence uses KM. The precise primary-source review and actual joint interface are recorded in \texttt{docs/KMTT\_INPUTS\_WORKING.md}.

\textbf{Derived OR: orbit transference.} The required real-weight specialization of Mousavi's related sequences-paper Proposition 1.4 is now proved by \texttt{orbitVNInput\_\allowbreak{}of\_\allowbreak{}finiteVNInput}. The complete supplied source is preserved unchanged at \texttt{manuscript/}\allowbreak\texttt{related-original/}\allowbreak\texttt{PSSequences.tex}; no main result of that paper is assumed. The proof derives \texttt{OrbitVNInput} from TV by explicit complex duality, exact spatial averaging over $N^d$ translates, degree-preserving polynomial shifts, supported orbit samples and integration. Essential bounds and finite-average integrability are proved. Its former OR-B bridge is discharged for this interface. The published finite-estimate supremum bridge TV-B remains external. The pointwise and new variation branches use this derived orbit estimate.

\textbf{MW: ordinary weighted pointwise convergence.} Mirek--Wan--Wright, arXiv:2505.15549v2, Theorem 1.1(ii), proves a.e. convergence for finite $1<q_i<\infty$ with $\sum_iq_i^{-1}\leq1$. \texttt{OrdinaryPrimePointwiseInput} uses bounded probability-space inputs and the shifted family $P_i(n-1)$. The finite-$L^q$ embedding and source-to-integer-action application are the unformalized translation MW-B; the polynomial-shift degree identities are independently proved. Only Theorem 1.2 depends on the MW pointwise interface. The additional MW-V interface quotes the same paper's Theorem 1.1(iv), equation (1.5), for $r>2$ and $\lambda>1$. It supplies an essentially $1$-bounded probability-space seminorm estimate uniform over the positive $\lambda$-lacunary sequence and system. Its separate external bridge MW-V-B specializes finite exponents $q_i=2\ell$, $q=2$, uses $L^\infty\subset L^{2\ell}$ and $L^2\subset L^1$ on a probability space, and passes from the published full variation norm to the seminorm. Polynomial shift/degree identities are already checked; the complete external-theorem application remains unformalized. Only new Corollary 4.2 depends on MW-V. The original two main theorems still jointly use six interfaces; the old standalone results use seven, and the extended project uses eight. These counts refer to interfaces, not independent published theorems.

\textbf{PNT: ordinary weighted mass.} \texttt{PrimeNumberTheoremInput} states $N^{-1}\sum_{n\leq N}\Lambda(n)\to1$. Tao's \emph{246B Notes 4: The Riemann zeta function and the prime number theorem}, 12 February 2021, Theorem 6, states this von Mangoldt form. Integer-scale specialization is included in the quoted input. This is not silently inferred from an unrelated prime-counting normalization. Higher-prime-power removal is proved internally. The optional no-Siegel footnote is not used; Green's Hexagon:2610.00311v2 Theorem 1.1 is not an additional project premise.

\else
\textbf{SS} is the directly published polynomial-phase consequence of Shao--Sun, arXiv:2605.18676v1, Theorem 2.1 and Definition 4.5. \textbf{TV} is Ter\"av\"ainen, arXiv:2401.03174v2, Theorem 4.1. The universal-phase/supremum bridge TV-B remains external. SS and TV give Proposition 2.3. TV also proves the orbit transference used by the pointwise branch.

\textbf{KM-U} is KMTT, arXiv:2608.19525v1, Proposition 5.2; \textbf{KM-R} is Lemma 5.3, with its small-scale alternative and coordinate-dependent resolutions. \textbf{KM-M} records (5.2) and (5.15). The actual \texttt{KMTTCompatibleInput} selects the same comparison and mass parameters jointly; \textbf{KM-J} is its explicitly unformalized fixed-cutoff and constant-reselection bridge. The standalone \texttt{KMTTUnweightedInput} is subsumed when this joint input is used. Only recurrence uses these predicates.

\textbf{OR is derived from TV}: the required real-weight specialization of the related sequences-paper Proposition 1.4 is proved internally, including shifts, support, essential bounds and integrability. \textbf{MW} is Mirek--Wan--Wright, arXiv:2505.15549v2, Theorem 1.1(ii), via \texttt{OrdinaryPrimePointwiseInput}; its source translation MW-B remains external. \textbf{PNT} is the ordinary von Mangoldt mass limit, directly stated in Tao's 246B Notes 4, Theorem 6. Theorem 1.2 has exactly SS, TV, PNT and MW as its substantial external premises. The new MW-V interface quotes MWW Theorem 1.1(iv), equation (1.5); its finite-$L^q$, H\"older and full-norm-to-seminorm translation is external MW-V-B. Corollary 4.2 uses exactly SS, TV and MW-V, without PNT. The original main-theorem premises are unchanged. The optional no-Siegel footnote is not used.
\fi

Primary sources were read in the cited versions. Provenance verification does not formalize an external theorem or discharge its source-formulation bridges. See the repository assumption catalogue and reference review for the actual interfaces.
